\section{Related Work}

\subsection{Linear Attention}

\textbf{Linear Transformer} (Katharopoulos et al., 2020) first replaced softmax 
with a kernel feature map. \textbf{Performer} (Choromanski et al., 2021) uses 
positive orthogonal random features. \textbf{cosFormer} (Qin et al., 2022) uses 
ReLU with cosine reweighting. \textbf{Nystr mformer} (Xiong et al., 2021) uses 
Nystr m approximation with landmarks.

\subsection{Landmark Methods}

\textbf{Nystr mformer} uses Nystr m landmarks selected via a downsampling rule. 
\textbf{Set Transformer} (Lee et al., 2019) uses inducing points. \textbf{Perceiver} 
(Jaegle et al., 2021) uses a learned latent array.

\textbf{Our distinction.} We select landmarks \textbf{per-example} via a learned 
scorer that operates on the RFF features. The selection is content-dependent 
and trained end-to-end.

\subsection{Local-Global Hybrid}

\textbf{Longformer} (Beltagy et al., 2020) combines local sliding-window 
attention with global tokens. \textbf{BigBird} (Zaheer et al., 2020) adds 
random attention. \textbf{Hymba} (Dong et al., 2024) and \textbf{Jamba} (Lieber 
et al., 2024) fuse attention and SSM.

\textbf{Our distinction.} We fuse \textbf{local sliding window} with **landmark 
linear attention**, both derived from the same Q/K/V projections.

\subsection{Softmax Variants}

\textbf{Softpick} (Gu \& Dao, 2024) uses ReLU-damped softmax for stable 
masked attention. \textbf{Sparsemax} (Martins \& Astudillo, 2016) projects 
onto the simplex. \textbf{Entmax} (Peters et al., 2019) uses $\alpha$-entmax.

\textbf{Our use.} We use softpick to handle the case where all rows of the 
attention matrix are masked (which causes standard softmax to give 
uniform attention).

\subsection{Random Fourier Features}

\textbf{Rahimi \& Recht (2007)} introduced RFF for kernel approximation. 
\textbf{Yu et al. (2016)} proposed orthogonal random features with lower 
variance. We extend RFF with \textbf{dual-scale} features.

\subsection{Position of Our Work}

\begin{table}[H]
\centering
\small
\begin{tabular}{lllll}
\toprule
\textbf{Method} & \textbf{Linear} & \textbf{Landmark} & \textbf{Local} & \textbf{Softmax} \\
\midrule
Linear Transformer & Yes & No & No & -- \\
Performer & Yes & No & No & RFF \\
Nystr mformer & Yes & Fixed & No & -- \\
Longformer & No & No & Yes & Softmax \\
\textbf{LEKAv2} & \textbf{Yes} & \textbf{Learned} & \textbf{Yes} & \textbf{Softpick} \\
\bottomrule
\end{tabular}
\end{table}



\subsection{Comparison Table}

\begin{lstlisting}

---
\end{lstlisting}